\documentclass{article}
\usepackage[utf8]{inputenc}
\usepackage[russian]{babel}
\usepackage{amsmath}
\usepackage{amssymb}
\usepackage{mathtools}
\usepackage{amsthm}
\usepackage[a4paper, left=2cm, right=2cm, top=2cm, bottom=2cm]{geometry}

\theoremstyle{plain}
\newtheorem{theorem}{Теорема}
\newtheorem{lemma}[theorem]{Лемма}
\newtheorem{corollary}[theorem]{Следствие}

\theoremstyle{definition}
\newtheorem{definition}{Определение}
\newtheorem{example}{Пример}

\theoremstyle{remark}
\newtheorem{remark}{Замечание}

\setlength{\parindent}{1.5em}

\title{Билет 11. Переход к пределу в несобственных интегралах первого рода, зависящих от параметра. Непрерывность интеграла по параметру.}
\author{}
\date{}

\begin{document}

\maketitle

\section*{Постановка задачи}

Рассматриваем функцию $f(x, y)$, заданную на декартовом произведении $[a, +\infty) \times Y$, где $Y$ -- произвольное числовое множество. То есть $(x, y) \in [a, +\infty) \times Y$.

Считаем, что $\forall \overline{y} \in Y$ существует несобственный интеграл первого рода
\[
\exists\int_{a}^{+\infty} f(x, \overline{y}) \, dx.
\]

\section*{Теорема о переходе к пределу под знаком интеграла}

\begin{theorem}[О переходе к пределу под знаком интеграла]
Пусть $y_0$ -- конечная или бесконечная предельная точка множества $Y$. Пусть выполнены следующие условия:
\begin{enumerate}
    \item Функция $f(x, y)$ равномерно на любом конечном отрезке $[a, A]$, $A > a$ при $y \to y_0$ сходится к предельной функции $\varphi(x)$, то есть
    \[
    f(x, y)\rightrightarrows^{[a, A]}_{y \to y_0} \varphi(x)
    \]
    \item Интеграл $\int_{a}^{+\infty} f(x, y)dx = G(y)$ сходится равномерно на $Y$;
\end{enumerate}
Тогда существует предел
\[
\lim_{y \to y_0} G(y) = \int_{a}^{+\infty} \varphi(x)dx.
\]
\end{theorem}

\begin{proof}[Доказательство]
\textbf{Существование.} Докажем, что интеграл $\int_{a}^{+\infty} \varphi(x)dx$ сходится.

Если $f(x, y)$ (Условие 1 теоремы) на любом конечном промежутке сходится к $\varphi(x)$, а сама функция $f(x, y)$ интегрируема на любом таком отрезке, то по Теореме об интегрируемости предельной функции:

\begin{theorem}Если $f(x, y)$ интегрируема $\forall \overline{y} \in [a, b]$, то и ее предельная функция интегрируема на этом отрезке.
\end{theorem}
мы можем утверждать, что $\varphi(x)$ интегрируема на любом отрезке $[a, A]$. \\
\textbf{Сходимость.} Используем критерий Коши сходимости для несобственных интегралов. Нужно доказать, что
\[
\forall \varepsilon > 0 \quad \exists B : \forall A', A'' > B \quad 
\left| \int_{A'}^{A''} \varphi(x)dx \right| < \varepsilon.
\]
Необходимо сформулировать правило как для заданного $\varepsilon$ можно найти $B$.

Берем $\varepsilon > 0, \frac{\varepsilon}{2} > 0.$

Рассматриваем $G(y) = \int_{a}^{+\infty} f(x, y)dx$ - несобственный интеграл 1 рода. По условию 1 теоремы, он равномерно сходится, значит для него выполнен критерий Коши:
\[
\frac{\varepsilon}{2} > 0 \quad \exists B : \forall A', A'' > B \quad 
\left| \int_{A'}^{A''} f(x, y)dx \right| < \frac{\varepsilon}{2}.
\]
Берем это значение $B$ для критерия Коши сходимости интеграла от $\varphi(x)$. Чтобы доказать выполнение критерия, необходимо доказать, что:
\[
\left| \int_{A'}^{A''} \varphi(x)dx \right| < \varepsilon.
\]
Используем неравенство треугольника:
\[
\left| \int_{A'}^{A''} \varphi(x)dx \right| \leq
\left| \int_{A'}^{A''} f(x, y) \, dx \right| +
\left| \int_{A'}^{A''} \varphi(x) - f(x, y) \, dx \right|.
\]
Первое слагаемое меньше $\frac{\varepsilon}{2}$ по критерию Коши для $f(x, y)$.

Для оценки второго слагаемого используем условие 1 из теоремы, $f(x, y)$ равномерно сходится к $\varphi(x)$ на любом отрезке $[a, A]$. Значит,
\[
\frac{\varepsilon}{2\left|A'' - A'\right|} > 0 \quad \exists \tilde{y} \in Y :
\left| f(x, \tilde{y}) - \varphi(x) \right| < \frac{\varepsilon}{2\left|A'' - A'\right|}.
\]
Тогда $\left|\varphi(x) - f(x, y)\right| < \frac{\varepsilon}{2\left|A'' - A'\right|}$, значит и вся сумма будет меньше $\varepsilon$. \\
\textbf{Сходимость.}
Доказываем, что \[\lim_{y \to y_0} G(y) = \int_{a}^{+\infty} \varphi(x)dx.\]
По определению предела функции нужно доказать, что
\[
\forall \varepsilon > 0 \quad \exists \delta > 0 : \forall y \in Y, \quad 0 < |y - y_0| < \delta \quad
\left| G(y) - \int_{a}^{+\infty} \varphi(x)dx \right| < \varepsilon.
\]
Рассматриваем этот модуль, испольузуем неравенство треугльника:
\begin{align*}
\Bigl| \int_{a}^{+\infty} & f(x, y)dx - \int_{a}^{+\infty} \varphi(x)dx \Bigr| \\
&\leq \Bigl| \int_{a}^{+\infty} f(x, y)dx - \int_{a}^{A} f(x, y)dx \Bigr| \\
&\quad + \Bigl| \int_{a}^{A} \bigl(f(x, y) - \varphi(x)\bigr) dx \Bigr| \\
&\quad + \Bigl| \int_{a}^{A} \varphi(x)dx - \int_{a}^{+\infty} \varphi(x)dx \Bigr|.
\end{align*}
Необходимо, чтобы сумма была меньше $\varepsilon$. Для этого необходимо,
чтобы каждое слагаемое было меньше $\frac{\varepsilon}{3}$.

Выбираем $A$. Рассматривая последний модуль, мы доказали, что интеграл $\int_{a}^{+\infty} \varphi(x)dx$ сходится. То есть при достаточно большом $A$ последний модуль будет меньше $\frac{\varepsilon}{3}$.
Выбирая то же $A$, первый модуль тоже будет меньше $\frac{\varepsilon}{3}$.

Второй модуль имеет конечные пределы. Его можно сделать меньше $\frac{\varepsilon}{3}$, используя условие 1 теоремы о равномерной сходимости $f(x, y)$ к $\varphi(x)$ на отрезке $[a, A]$. Для этого нужно выбрать такое $\delta > 0$, что при $0 < |y - y_0| < \delta$ будет выполнено требуемое неравенство.

Таким образом, все три слагаемых будут меньше $\frac{\varepsilon}{3}$, значит сумма будет меньше $\varepsilon$. Теорема доказана.
\end{proof}

\end{document}